% D3 — maquette d'une séance du livre du formateur.
% Contenu : bir/formateur/seance-08.md.
\documentclass{BIRdoc}
\BIRdocument{Livre du formateur}{Séance 8 --- Le spectre et les bandes}
\begin{document}

\BIRtitre{Séance 8 --- Le spectre et les bandes}{Domaine C · 1 h 30 · chapitre 8 du manuel · fiche d'activité nº 8}

\begin{BIRencadre}{Ce que l'élève doit emporter}
\begin{itemize}
\item Le spectre radio va de 30 kHz à 300 GHz ; il se découpe en domaines, dont HF, VHF et UHF.
\item Les fréquences sont partagées entre de nombreux usages ; ce partage est organisé entre les pays, puis dans chaque pays.
\item Les radioamateurs disposent de bandes, nommées par leur longueur d'onde.
\item Un spectre se lit en fréquence et en force ; une chute d'eau ajoute le temps.
\end{itemize}
\end{BIRencadre}

\section*{Avant la séance}

\textbf{Acquis nécessaires :} fréquence et longueur d'onde (séance 7).

\begin{minipage}[t]{.48\linewidth}
\textbf{Par groupe de 2 ou 3}
\begin{itemize}
\item un ordinateur avec un logiciel de réception SDR installé et essayé ;
\item une clé SDR et son antenne ;
\item la fiche élève nº 8, un crayon.
\end{itemize}
\end{minipage}\hfill
\begin{minipage}[t]{.48\linewidth}
\textbf{Pour l'encadrant}
\begin{itemize}
\item un poste ou un récepteur relié au vidéoprojecteur ;
\item un WebSDR en secours, et pour montrer les bandes HF que la clé ne reçoit pas.
\end{itemize}
\end{minipage}

\medskip
\textbf{À préparer.} Relever à l'avance, depuis la salle, quatre ou cinq signaux
sûrs : deux stations FM, un multiplex DAB+, le relais radioamateur local, une
balise. Noter leurs fréquences : ce sont les points de repère de la séance.

\section*{Déroulé}

\begin{tabularx}{\linewidth}{P{14mm}P{30mm}LL}
\BIRentete \BIRth{Durée} & \BIRth{Phase} & \BIRth{L'encadrant} & \BIRth{Les élèves} \\
5 min & accroche & projette une chute d'eau de la bande FM : « combien de stations ici ? » & comptent les bandes verticales \\
10 min & le spectre & déroule le tableau des domaines ; place la radio FM, le Wi-Fi, le relais local & situent trois fréquences dans HF, VHF ou UHF \\
10 min & le partage et les bandes & pourquoi partager ; qui décide ; six bandes radioamateur & retrouvent le nom d'une bande par le calcul de la séance 7 \\
\rowcolor{BIRclair} 50 min & manipulation & circule, dépanne, relance & dressent leur carte du spectre (fiche nº 8) \\
15 min & bilan & fait comparer les cartes ; corrige trois questions & complètent « à retenir » ; répondent aux questions \\
\bottomrule
\end{tabularx}

\section*{Conduite de la manipulation}

\begin{enumerate}
\item \textbf{Prise en main (10 min).} Tout le monde sur une station FM connue.
  Faire repérer l'axe des fréquences, le pic, la trace dans la chute d'eau.
  Couper le son, puis le remettre : la trace ne dépend pas du son.
\item \textbf{Exploration guidée (25 min).} Les groupes cherchent les signaux de
  la fiche, dans l'ordre. Pour chacun : fréquence, domaine, aspect à l'écran.
\item \textbf{Exploration libre (15 min).} Chaque groupe trouve un signal qui
  n'est pas sur la fiche et le décrit aux autres.
\end{enumerate}

\begin{BIRalerte}{Ce qu'on écoute, ce qu'on regarde}
Les élèves n'écoutent que la radiodiffusion et les bandes radioamateur. Tout
autre signal est observé à l'écran, son coupé : on note sa fréquence et sa
forme, pas son contenu.
\end{BIRalerte}

\section*{Ce qui coince souvent}

\begin{description}[leftmargin=0pt,itemsep=3pt]
\item[Confondre les deux axes.] « Plus à droite » est lu comme « plus fort ».
  Faire montrer du doigt : droite-gauche, c'est la fréquence.
\item[Confondre spectre et oscillogramme.] À la séance 7, l'axe horizontal
  était le temps. Ici, c'est la fréquence. Le dire explicitement.
\item[Croire qu'une bande est une seule fréquence.] Faire compter les stations
  visibles dans la bande FM : une bande, beaucoup de fréquences.
\item[Croire que « 2 m » est une distance ou une taille d'antenne.] Refaire le
  calcul : 300 divisé par 145.
\item[Prendre un parasite pour un émetteur.] Un pic qui reste quand on débranche
  l'antenne vient de l'ordinateur ou de la clé.
\end{description}

\section*{Pour aller plus loin}

\begin{description}[leftmargin=0pt,itemsep=3pt]
\item[Statut primaire ou secondaire.] Sur certaines bandes, les radioamateurs
  sont prioritaires ; sur d'autres, ils ne doivent pas gêner le service
  principal. La bande des 2 m est à statut primaire, celle de 430 à 434 MHz
  à statut secondaire.
\item[Les trois régions de l'UIT.] Le monde est découpé en trois régions ; les
  bandes n'y sont pas toutes identiques. La France métropolitaine est en région 1.
\item[Les plans de bande.] À l'intérieur d'une bande, les radioamateurs se
  répartissent eux-mêmes les usages : télégraphie ici, voix là, balises ailleurs.
\item[L'échelle verticale du spectre] est graduée en décibels : chaque
  graduation correspond à un rapport, pas à un écart.
\end{description}

\section*{Corrigé de la fiche}

Les fréquences relevées dépendent du lieu. Sont attendus : le bon domaine pour
chaque signal, une description juste de la trace, et la longueur d'onde
calculée à 10 \% près.

\begin{tabularx}{\linewidth}{P{45mm}P{30mm}L}
\BIRentete \BIRth{Signal} & \BIRth{Domaine} & \BIRth{Aspect attendu à la chute d'eau} \\
station FM & VHF & bande large, continue \\
multiplex DAB+ & VHF & bloc très large, uniforme, à bords nets \\
relais radioamateur & VHF ou UHF & bande étroite, par moments seulement \\
balise & selon la bande & trait fin, continu ou découpé \\
\bottomrule
\end{tabularx}

\medskip
\textbf{Questions de la séance :} RC0801 à RC0810.

\end{document}
